\section{A fixed marked section and its calibration budgets}
\label{sec:ex-marked-portfolio-calibration}

The marked finite-volume route needs a single local section at each label
degree. That section must be chosen before the volume, cutoff, or derivative
budget is known. It must also price every later derivative request at that
degree. We construct that section first, then show how the schedule selects
its budgets and how the approximate field change is made invertible.

\subsection{One section for every later derivative budget}

Let $H$ be the fixed holomorphic gauge action. At a fixed degree $d$, let
$E_d$ be a finite-dimensional complex label space with a real homogeneous
basis $(e_j)_{j\in J_d}$. The basis contains the unit and has grades
$(r_j,s_j)$ with $r_j\leq d$. The algebra used here is not filtered by raw
link-entry count. Let $\mathfrak S$ be the collection of based and
transported curvature slots on the cover lattice, let $I$ be the real
color-coordinate basis, and put $\mathrm{Chip}=\mathfrak S\times I$.
For $c=(S,a)\in\mathrm{Chip}$ let $x_c(U)$ be color coordinate $a$ of
curvature slot $S$ at configuration $U$. Define $\mathcal P_{\geq s}$ as
the complex span of the finite products $\prod_{\nu=1}^m x_{c_\nu}$ with
$m\geq s$. Thus one factor means one curvature-coordinate slot, even
though it is itself polynomial in link entries. This is the source
filtration called $\mathrm{Powered}$.

The input to the section construction is an imprint: linear maps
$A_d:E_d\to E_d$ and $B_d:E_d\to\mathcal P_{\geq0}$ with unit normalization
$A_d1=1$, $B_d1=1$, gauge invariance of $B_dp$,
$\operatorname{Forms}(H,d,B_dp,A_dp)$, and $B_dp\in\operatorname{Micro}$
whenever $p$ belongs to the equation ideal, together with the campaign's
exact-action identity. They also have
the following filtration clauses. The source stamp of a grade is
\[
                         \operatorname{st}(r,s)=(s,r),
\]
and stamps are compared lexicographically in the displayed order. Say
that a label has stamp at least $\operatorname{st}(r,s)$ (respectively
strictly above it) if each of its nonzero homogeneous components, of
grade $(r',s')$ say, has
$\operatorname{st}(r',s')\geq\operatorname{st}(r,s)$ (respectively $>$).
If $p$ has stamp at least $\operatorname{st}(r,s)$, then $A_dp-p$ has
stamp strictly above $\operatorname{st}(r,s)$, and
$B_dp\in\mathcal P_{\geq s}$. By linearity it is equivalent to require
these two clauses for homogeneous $p$ of grade $(r,s)$ only: a component
of grade $(r',s')$ with larger stamp has $s'\geq s$, because increasing
the stamp never lowers the factor coordinate, and
$\mathcal P_{\geq s'}\subseteq\mathcal P_{\geq s}$. In addition, $A_d$ and
$B_d$ each commute with
the signed-coordinate action and the real structure; equivariance of only
$B_d$ would not suffice to pass it through $A_d^{-1}$. A campaign without
these two equivariances is turned into such an imprint by the finite group
and conjugation averages recalled next; the campaign producer itself is a
separate input to this local theorem. It is constructed in
Section~\ref{sec:marked-campaign-producer}.

For clarity, the symmetrization just mentioned is finite-dimensional. If
$\Phi_s$ is the joint signed action on labels and fields, and $\mathfrak c$
is the joint real-conjugation involution, put
\[
 F_0=\frac1{|\mathcal S|}\sum_{s\in\mathcal S}
       \Phi_s\,F\,\theta_{s^{-1}},\qquad
 F_1=\frac12\bigl(F_0+\mathfrak c F_0\mathfrak c\bigr),
\]
where $\mathcal S$ is the signed-coordinate group. Each summand is still a
campaign: the grade is preserved, the correction remains strictly above
its source stamp, the field remains in $\mathcal P_{\geq s}$, and the unit
and exact action identities are preserved. A finite average preserves
these linear conditions. Reindexing the first sum by left multiplication
shows that $F_0$ commutes with every signed action; the involution identity
$\mathfrak c^2=1$ shows $F_1$ commutes with conjugation, and the two
actions commute, so both equivariances hold for $F_1$. Both averages fix
the unit. This is the mathematical content of passing from a campaign to
an imprint; the existence of the unsymmetrized campaign from the action is
not reproved in this section.

\begin{theorem}[Fixed-stencil marked portfolio from an imprint]
\label{thm:ex-marked-portfolio-section}
Given an imprint on $E_d$, there are a finite centered stencil $S_d$, a linear section
\[
       T_d:E_d\longrightarrow\mathbb C[\text{link entries on }S_d],
\]
and an invertible linear regrading $R_d:E_d\to E_d$ such that
\begin{enumerate}
\item $0\in S_d$ and $S_d$ is invariant under signed coordinate
  permutations;
\item evaluating $T_d$ at the coordinates of $R_dp$ on every link field
  gives $B_dp$ exactly;
\item $T_d$ is gauge invariant, has the correct unit, is compatible with
  signed pullbacks and conjugation, and has the prescribed action
  projection;
\item each basis entry has a finite pose supported in $S_d$ whose degree
  is at least $s_j$;
\item for every derivative order $k$ and power $P$, there is a finite
  number $C_{d,k,P}$ such that each basis entry has source rate
  $C_{d,k,P}$ at factor degree $s_j$. This is a bound on the healthy part
  of the translated stencil, for every budget point of order $k$ and
  power $P$, torus, anchor, and tested direction list allowed by that
  budget point;
\item one may take $R_d=A_d$. Then $R_d$ commutes with the
  signed-coordinate action and with the real structure, and for every
  $p$ in the equation ideal the field obtained by evaluating $T_d$ at the
  coordinates of $R_dp$ belongs to $\operatorname{Micro}$.
\end{enumerate}
The section and stencil depend on $d$ and the fixed action/frame only;
they do not depend on $k,P$, torus size, or cutoff.
\end{theorem}
\begin{proof}
First invert the triangular regrading. On the finite grade decomposition,
the condition that $A_dp-p$ has strictly higher grade makes $A_d$ an
upper-triangular block matrix with identity diagonal. Back substitution
solves $A_dx=y$ grade by grade: at the lowest grade, $x$ has the same
component as $y$; after subtracting that component, the next grade is
solved in the same way. Finiteness of the grade set makes this process
terminate and gives a unique $x$. Thus $R_d=A_d$ is invertible, and the
same induction shows that its inverse preserves the lower grade bound.

For each basis vector set $f_j=B_d(R_d^{-1}e_j)$. By the filtration
property, $f_j$ is a finite linear combination of monomials, each with at
least $s_j$ curvature-coordinate factors. Each monomial uses only a
finite list of slots. A product is represented by concatenating its two
lists of slots. A sum is represented
by taking the disjoint union of the lists and adding the corresponding
polynomials. The empty list represents constants. Induction over the
finite monomial expansion therefore gives, for each $j$, a finite slot
list and a pose polynomial of degree at least $s_j$ realizing $f_j$.

There are only finitely many basis vectors. The union of their slot
stencils is finite, so choose a centered box containing it. Enlarge the
box symmetrically if needed; this preserves containment and makes the
stencil invariant under signed coordinate permutations. On this one box,
the polynomial realization is constructed for every basis vector. Extend
these values linearly from the basis to define $T_d$. Linearity gives the
exact evaluation identity for every $p$, not merely for basis labels.

The evaluation identities also transfer all structure. Gauge invariance,
the unit normalization, signed covariance, conjugation reality, and the
action-projection condition are the corresponding properties of $B_d$.
Because evaluation agrees on the entire basis, the section has these
properties on all labels. The basis poses give the common-support and
degree witnesses. The last item of the theorem is immediate from the
choice $R_d=A_d$: the two equivariances of $A_d$ are imprint hypotheses,
and by the exact evaluation identity the field attached to $R_dp$ is
$B_dp$, which is microscopic for $p$ in the equation ideal.

It remains to show the derivative budgets for all later queries. The
source rate statement is localized: for a stencil $S$, order $k$, power
$P$, side $n$, and anchor $o$, it controls derivatives after expansion
at $o$ on $\mathrm{Healthy}(\theta,S+o)$, for direction lists obeying
the budget point's direction bound at $o$. The extra derivative-rate
parameter here is zero, so the threshold factor remains exactly
$\theta^s$ at every tested derivative order.
Each color-coordinate curvature slot has rate degree one: on
configurations healthy on the translated stencil, the slot and each of
its tested derivatives through order $k$ have norm at most a constant
times $\theta$, uniformly in the torus and the anchor. This is the slot
estimate proved in Section~\ref{sec:local-sections} (the value is bounded
by health and the isometry of parallel transport, and every derivative of
positive order $\ell$ by a constant times $\theta^\ell\leq\theta$ because
each tested direction on a link of the stencil has size at most a
constant times $\theta$). A real color coordinate is a bounded linear
functional of the slot, so it obeys the same bound. For a binary product,
the product rule gives at most $2^k$ terms at order $k$. Iterating over a
product of $m$ factors gives at most $m^k$ terms; in each pose the number
$m$ is fixed, so this changes only its
finite constant. Factor degrees add under multiplication. Hence a monomial of
at least $s_j$ chips has rate degree at least
$s_j$. Finite sums add their constants, and lowering a rate from a larger
degree to $s_j$ is valid because $0<\theta\leq1$. Applying this to the
finite pose expansion of each basis field gives one finite constant
$c_j(k,P)$ for the exact localized rate above. Summing basis estimates
with coefficient moduli gives the linear price for arbitrary labels. No
uniformity in $d,k$, or $P$ is claimed.
\end{proof}

Item (5) is the currency furnished by the finite section. In the
application, the profile at stage $i$ requests the particular order
$k_i$ and power $P_i$ appearing in that profile. For each $d\leq i$, choose
one available amount $a_{di}\geq0$ for this request. It remains a valid
credit after replacing a negative amount by zero, since all budgets are
nonnegative. The section is not rebuilt when these requests change.

\subsection{Buying all stage and degree conditions}

For a stage $i$, let $c_i$ be its column count, let $q_d$ be the stencil
size of the chosen section together with the fixed local section, and let
$\alpha_i$ be its support exponent. The source sphere price is
\[
 \mathsf S_i(q_d,a)=2+e^{81\alpha_i}(a+c_i+2q_d).
\]
For each fixed pair $d\leq i$, the chosen credit $a_{di}$, the sphere price,
and the coefficient $2(1+c_i)$ are finite. The orbit cost of the
degree-$d$ frame,
\[
 \mathcal O_d=1+\sum_{g}\sum_{j\in J_d}\|\rho_ge_j\|_1,
\]
where $g$ runs over the finite list of signed coordinate permutations
with or without conjugation and $\rho_g$ is its action on coefficient
vectors, is finite and at least one. Hence its sliver
\[
       \delta_d=\frac{1}{2\,\mathcal O_d^2}
\]
is strictly positive. Since the volume purse is
$\mathsf P(L)=L^{1/8}\to\infty$, choose an integer $n_{di}$ such that
\[
 2(1+c_i)\mathsf S_i(q_d,a_{di})
       \leq \mathsf P(L)\delta_d\qquad(L\geq n_{di}).
\]
The credit clause and nonnegativity hold at every such $L$. Separately,
for each $j$ choose $v_j$ so that
\[
       (1+2c_j)\,\mathsf D(L)\leq\overline{\mathsf P}(L)
                  \qquad(L\geq v_j),
\]
where the source door toll and purse are
$\mathsf D(L)=\mathrm{rate}(L)=L^{-2/5}$ and
$\overline{\mathsf P}(L)=\mathrm{barunit}(L)=L^{-3/8}$. Their ratio is
$L^{-1/40}$, so the inequality holds once $L^{1/40}\geq1+2c_j$. For a fixed stage $i$, take
the maximum of the finitely many $n_{di}$ with $d\leq i$ and $v_j$ with
$j\leq i$. This single threshold $n_i$ proves all clauses required at
stage $i$ simultaneously. The resulting nonnegative table $a_{di}$ and
thresholds $n_i$ are exactly the same section's all-stage purchase data.

The later delayed budget has a similar finite-threshold argument. Write
$\beta_F(x,L)=x+b_L$, where $b_L$ is the fixed logarithmic scale shift,
and write $\mathsf f_d(x)=x^{-w_d}$ for the reciprocal tail. For the
public schedule $w_d=\max_js_j+4$, the maximal factor grade of the
degree-$d$ frame plus four; the following delay argument works for any
fixed positive integer exponent. The tail
comparison needed by the delay is, with $y=\beta_F(x,L)=x+b_L$,
\[
       y^{-1}L^{d+1}\,\mathsf f_d(x)\leq\mathsf f_d(y).
\]
For fixed $L,d$, and all sufficiently large $x$, one has $x\geq1$,
$b_L\leq x$, and $y\geq\max\{1,L^{d+1}2^{w_d}\}$; the last condition
holds eventually because $y\to\infty$ with $x$. In particular
$0<y\leq2x$. After multiplying by the positive denominators the tail
comparison is equivalent to
\[
       L^{d+1} y^{w_d}\leq yx^{w_d}.
\]
The left side is at most $L^{d+1}2^{w_d}x^{w_d}$, which is at most the
right side because $y\geq L^{d+1}2^{w_d}$. No sign condition on the
shift $b_L$ is used. The remaining conditions at a fixed
stage-profile pair are just that $x$ exceed the fixed profile cost and
the finitely many credits requested at degrees $d$ below that profile.
They therefore hold eventually as well. For a fixed stage, only that
stage and its adjacent predecessor are tested; the conjunction is finite.
Choose a threshold for every stage and replace these thresholds by their
prefix maxima, also taking the maximum with the original entrance delay.
The new entrance delay is monotone, dominates the old one, and satisfies
every adjacent-stage inequality whenever the actual delay exceeds it.
This proves the delay calibration rather than assuming a single threshold
works for all stages at once.

At the selected history caller, these are the three eventual conditions
needed by the marked budget package. For each fixed degree $d$, the section
credit is queried at the profile of the actual stage $J(N)$. The strict
reference-depth margin implies $d\leq J(N)$ for all sufficiently large
$N$. The purchased stage table supplies the credit and flight inequality
at that stage, evaluated at its preceding block side $L_{N-1}$; this side
is at least the stage threshold $n_{J(N)}$ provided the thresholds $n_i$
do not exceed the floor length of stage $i$, which bounds from below
every block side used at that stage. This comparison of thresholds with
floor lengths is an explicit condition on the schedule of the caller; it
is not proved in this section. The
delay purchase supplies the allowance comparison. At level $N$ the
selected coupling $\widehat\beta_N$ is at least the entrance delay of the
stage $J(N+1)$, and $J(N)\leq J(N+1)\leq J(N)+1$, so the stage $J(N+1)$
and the profile $J(N)$ form one of the adjacent pairs tested by the delay.
Its credit clause at $x=\widehat\beta_N$ gives
$a_{d,J(N)}\leq\widehat\beta_N$ for every $d\leq J(N)$. The credit
$b_N=a_{d,J(N)}$ is thus nonnegative and no larger than the running
allowance $\widehat\beta_N$. The tail clause of the same delay is used
separately: with the schedule identity
\[
 \widehat\beta_{i+1}=\beta_F(\widehat\beta_i,L_i)
       =\widehat\beta_i+c_G\log L_i,
\]
where $c_G$ is the slope constant of the gauge group, it reads
\begin{equation}\label{eq:ex-marked-height-jump}
 \widehat\beta_{i+1}^{-1}L_i^{d+1}\,\mathsf f_d(\widehat\beta_i)
      \leq\mathsf f_d(\widehat\beta_{i+1})\qquad(d\leq J(i)).
\end{equation}
This is the height recurrence
$\phi_iL_i^{d+1}\zeta_{d,i}\leq\zeta_{d,i+1}$ consumed in
Section~\ref{sec:actual-marked-portfolio-calibration}, for the public
choices $\phi_i=\widehat\beta_{i+1}^{-1}$ and
$\zeta_{d,i}=\mathsf f_d(\widehat\beta_i)$.
Thus, at every sufficiently late stage for this fixed degree,
there are simultaneously a valid derivative credit, the allowance bound,
and the sphere/flight bound. These are exactly the three fields of the
marked budget package used by the selected field-calibration call. It uses
the same portfolio section throughout: the delay changes available budget,
not the section or its stencil.

\subsection{A small coordinate error gives an invertible calibration}

The section is now fixed. At a given cutoff, the field calculation
produces, for each coefficient basis vector $e_j$, a proposed correction
$u_j$ satisfying
\[
        \|e_j-u_j\|_1\leq\delta_d,
\]
where $\|x\|_1$ is the sum of the absolute values of the finitely many
coordinates and $\delta_d=(2\mathcal O_d^2)^{-1}$. The required linear
calibration is not inferred from columnwise closeness alone; the factor
$\mathcal O_d$ accounts for the finite symmetry and the change from the
coordinate basis to the section's coefficient coordinates.

\begin{lemma}[Robust finite-dimensional calibration]
\label{lem:ex-marked-calibration-inverse}
Let $E$ be the coefficient space of the chosen frame, with its $\ell^1$
norm. Let $\mathcal O\geq1$ bound both the sum of the chip coefficients
and the $\ell^1$ norm of each symmetry action on the frame. For columns
$u_j$, define the source repainting map by
\[
 D x=\sum_{(g,j)}\chi_{g,j}(x)\,\rho_g u_j,
 \qquad
 x=\sum_{(g,j)}\chi_{g,j}(x)\,\rho_g e_j,
\]
where $g$ ranges over the finite signed/conjugation symmetry list,
$\rho_g$ is its action on the coefficient frame, and the fixed complex-
linear chip coefficients obey
$\sum_{g,j}|\chi_{g,j}(x)|\leq\mathcal O\|x\|_1$.
If $\|e_j-u_j\|_1\leq(2\mathcal O^2)^{-1}$ for every basis vector,
then $D$ is invertible and $\|D^{-1}\|_{1\to1}\leq2$.
\end{lemma}
\begin{proof}
Subtracting the two repainting formulas gives
$(I-D)x=\sum_{g,j}\chi_{g,j}(x)\rho_g(e_j-u_j)$.
The sum of chip moduli is at most $\mathcal O\|x\|_1$, and applying
the symmetry to a column error costs at most another factor $\mathcal O$.
Consequently
\[
       \|(I-D)x\|_1\leq\mathcal O^2(2\mathcal O^2)^{-1}\|x\|_1
                         =\tfrac12\|x\|_1.
\]
The triangle inequality gives $\|Dx\|_1\geq\frac12\|x\|_1$;
therefore $D$ is injective. Since $E$ is finite dimensional, it is
surjective and invertible. The same lower bound with $x=D^{-1}y$ gives
$\|D^{-1}y\|_1\leq2\|y\|_1$.
\end{proof}

\subsection{Application to the selected field calibration}

We now state how the fixed section is compared, at each cutoff, with the
local section of a selected family. The statement assembles two results
proved later, the one-body invitation of
Proposition~\ref{prop:ex-marked-invitation} and the same-history
calibration of Proposition~\ref{prop:actual-marked-figment}; neither of
them uses the present subsection.

Fix a selected family $C$ which makes the action, with its history, and a
degree $d$. Write $F$ for the grade-$d$ coefficient frame with basis
$(e_j)_{j\in J_d}$, $T$ for the local section of the family at degree
$d$, and $V=V_d$ for the fixed section of
Theorem~\ref{thm:ex-marked-portfolio-section}, expressed in this frame and
sound for the action of the family. At level $N$ write $a_N$ for the
mesh, $\theta_N$ for the budget threshold, $c_N$ for the column count,
$\alpha_N$ for the support exponent, $\sigma_N$ for the running scale,
$L_{N-1}$ for the preceding block side, and
$\mathsf S_N(q,b)=2+e^{81\alpha_N}(b+c_N+2q)$ for the sphere price. For a
section $W\in\{V,T\}$ and a coefficient vector $v\in F$, the normalized
origin field at $(N,p)$ is
\[
 \mathcal O^{W}_{N,p}(v)(U)=a_N^{-d}\,
     W\bigl(\theta_N^{-1}\mathbin{\cdot_F}v\bigr)(0,U),
\]
where $\cdot_F$ is the graded dilation multiplying the coordinate $j$ by
$\theta_N^{-s_j}$. The inputs are the following.
\begin{enumerate}
\item \emph{The eventual budget triple.} There is a sequence $b_N$ such
  that, for all sufficiently large $N$: the section $V$ has the credit of
  item~(5) of Theorem~\ref{thm:ex-marked-portfolio-section} with constant
  $b_N$ at the derivative order and power of the level-$N$ budget;
  $0\leq b_N\leq\widehat\beta_N$; and
  \[
   2(1+c_N)\mathsf S_N(q_d,b_N)\leq\mathsf P(L_{N-1})\delta_d .
  \]
  For a schedule satisfying the floor and delay conditions stated in the
  preceding subsection, the triple follows with $b_N=a_{d,J(N)}$; the
  discharge of those conditions on the selected public run is an open
  input.
\item \emph{The paid selected-history record}
  of Section~\ref{sec:actual-marked-portfolio-calibration}, with its
  allowances $u_N>0$ and heights $\zeta_{d,N}\geq0$, and a terminal Keeper
  at every $(N,p)$ (Section~\ref{sec:terminal-marked-construction}).
\item \emph{Slow growth of the reserves.} Call a real sequence $x_N$
  \emph{slow} if for every $\varepsilon>0$ one has
  $|x_N|\leq e^{\varepsilon\sigma_N}$ for all sufficiently large $N$
  along the extraction order under consideration. The sequences $u_N^{-1}$
  and $\zeta_{d,N}$ are slow for every extraction order.
\item The gauge group is simple, and the history has the opening schedule
  used for the centered root tests.
\end{enumerate}
Put
\[
 \mathsf Z_N=\bigl(\mathsf C(V)+2\mathsf C(T)\bigr)
                  \operatorname{Seat}_{d,N},
 \qquad
 \operatorname{Seat}_{d,N}=\sum_{j\in J_d}\theta_N^{-s_j},
\]
where $\mathsf C(W)\geq1$ denotes a fixed value-norm constant of the
section $W$: $|W(v)(o,U)|\leq\mathsf C(W)\|v\|_1$ for all coefficient
vectors, sides, anchors and configurations. Such a constant exists
because the finitely many basis columns of $W$ are continuous functions
on a compact configuration space. It is not a derivative price, and
$\operatorname{Seat}_{d,N}$ is not a number of sites: it bounds the
$\ell^1$ operator norm of the dilation $\theta_N^{-1}\cdot_F$. The
price attached to each constructed column is the explicit schedule amount
\[
 E_N=4(1+c_N)\mathsf S_N(q_d,b_N)
    +e^{81(\alpha_N+2)}\mathsf Z_N\,\zeta_{d,N}\,u_N^{-1}.
\]
The first summand pays the local stencil and sphere terms; the second
pays the forced correction after the reserve factor $\zeta_{d,N}/u_N$.

\begin{proposition}[Eventual invertible calibration of the fixed section]
\label{prop:ex-marked-solaris}
Under the four inputs above, for every sufficiently large $N$ and every
period $p$ there are columns $u_j=u_j(N,p)\in F$, $j\in J_d$, such that:
\begin{enumerate}
\item $\|e_j-u_j\|_1\leq\delta_d$ for every $j$. Consequently the
  repainting map $D_{N,p}$ built from these columns in
  Lemma~\ref{lem:ex-marked-calibration-inverse} is invertible and
  $\|D_{N,p}^{-1}\|_{1\to1}\leq2$.
\item For every $j$, every $k$ with $\operatorname{ref}(d)\leq k<N$,
  every $m$ opened at $k$, and every motion $s$ (a signed
  coordinate permutation, possibly composed with complex conjugation),
  the centered root test at depth $m$ for the finite-volume law at
  $(N,p)$, in the notation of
  Section~\ref{sec:actual-marked-portfolio-calibration}, satisfies
  \[
   \operatorname{rootTest}_{m;N,p}\Bigl(
      s\cdot\bigl(\mathcal O^{V}_{N,p}(e_j)
                  -\mathcal O^{T}_{N,p}(u_j)\bigr)\Bigr)
     \leq\operatorname{cider}(d,m,k)\,E_N\,
          e^{-\gamma\sigma_N},
  \]
  with the constant $\operatorname{cider}(d,m,k)$ and the exponent
  $\gamma=1/8$ of Proposition~\ref{prop:actual-marked-figment}. This
  assertion inherits the condition of that proposition: it depends on
  the served-representation step which is not proved there.
\item The sequence $E_N$ is slow for every extraction order.
\end{enumerate}
The threshold in $N$ depends on $d$ but not on $p$.
\end{proposition}
\begin{proof}
Fix $N$ so large that the following hold; each holds for all sufficiently
large $N$, uniformly in $p$. The three clauses of the budget triple hold
at $N$. One has $N>\operatorname{ref}(d)$. Every point $z$ of the union
$X$ of the stencils of $V$ and $T$ satisfies
$20(|z_i|+1)\leq\operatorname{side}(N,p)$ in each coordinate, because the
sides tend to infinity uniformly in $p$. Finally the three local reserve
conditions of Lemma~\ref{lem:ex-marked-sanded} and
\eqref{eq:ex-marked-cover} hold for the stencil $X$ and the two constants
$\mathsf C(V),\mathsf C(T)$; here simplicity of the gauge group is used,
through $\widehat\beta_N\to\infty$.

Fix $p$ and $j$. These are exactly the inputs of
Proposition~\ref{prop:ex-marked-invitation}, with the terminal Keeper,
$u_N>0$ and $\zeta_{d,N}\geq0$ from the second input. It supplies a
terminal coefficient $Z_j$ with leading coordinate $Z_j.x=e_j$, charge
bound $Q_N=\mathsf S_N(q_d,b_N)$, budget $E_N$, and field
$\mathcal O^{V}_{N,p}(e_j)$; these are the Cream data. Since
$\operatorname{ref}(d)<N$, Proposition~\ref{prop:actual-marked-figment}
applies on the same family and history. It gives one label $v_j$ with
\[
 \mathsf P(L_{N-1})\,\|e_j-X_N^{-1}v_j\|_1\leq2(1+c_N)Q_N
\]
and the displayed root-test bounds for
$\mathcal O^{V}_{N,p}(e_j)-C.O(v_j)$, where $X_N$ is the terminal
calibration isomorphism of the family. Put $u_j=X_N^{-1}v_j$. By the
definition of the family's observable, $C.O(v_j)=\mathcal
O^{T}_{N,p}(X_N^{-1}v_j)=\mathcal O^{T}_{N,p}(u_j)$; this proves the
second assertion. The flight inequality of the budget triple bounds the
right side of the last display by $\mathsf P(L_{N-1})\delta_d$. Dividing
by $\mathsf P(L_{N-1})>0$ gives $\|e_j-u_j\|_1\leq\delta_d$, and
Lemma~\ref{lem:ex-marked-calibration-inverse} gives the invertibility and
the inverse bound. The coefficient error is thus removed by one
invertible linear regrading, shared by all labels at this cutoff.

For the slow-growth assertion, all the components of $E_N$ are needed.
The class of slow sequences is closed under finite sums and products,
contains the constants, and contains every sequence dominated eventually
by a slow sequence. The credit satisfies $0\leq b_N\leq\widehat\beta_N$
eventually, and $\widehat\beta_N$ is slow because it is affine in
$\sigma_N$. The column count $c_N$ is slow by the accumulated
column-logarithm estimate. The opening schedule gives
$\alpha_N=o(\sigma_N)$, so $e^{81\alpha_N}$ and
$e^{81(\alpha_N+2)}=e^{162}e^{81\alpha_N}$ are slow. Hence
$\mathsf S_N(q_d,b_N)$ is slow. The inverse threshold $\theta_N^{-1}$ is
slow, being bounded by $\widehat\beta_N$; therefore the finite sum
$\operatorname{Seat}_{d,N}$ of its powers and $\mathsf Z_N$ are slow.
Finally $\zeta_{d,N}$ and $u_N^{-1}$ are slow by the third input. Each
summand of $E_N$ is a finite product of slow sequences. Omitting either
of the two reserve bounds of the third input would not justify this
conclusion.
\end{proof}

The section theorem and stage/delay arguments above prove the fixed-stencil
portfolio, all later credit requests, and, for a schedule satisfying the
stated floor and delay conditions, the three eventual field-budget
inputs. The local invitation and the same-history calibration
used in Proposition~\ref{prop:ex-marked-solaris} remain separate
substantive suppliers; they do not follow from existence of the section or
from the inverse lemma. Their source declarations, actual selected caller,
and remaining estimate boundary are listed in the evidence file.
